\section*{Bab \ref{ch:b2:oceanic-islands} --- Pulau Samudra}

\begin{solution}{exo:b2:oceanic-islands:1}
(a) Tahap praperisai, seperti Loihi kini. (b) \omterm{def:b2:oceanic-islands:stages}{Tahap perisai}. (c) Tahap pascaperisai. (d) Tahap
peremajaan. (e) Tidak ada lagi tahap erupsi: gunung api telah tenggelam di bawah terumbunya, yang kini menjadi \omterm{def:b2:oceanic-islands:reefs}{atol}.
\end{solution}

\begin{solution}{exo:b2:oceanic-islands:2}
(a) $2500 + 350\sqrt9 = \qty{3550}{m}$. (b) $350\times(\sqrt{25} - \sqrt9) = 350\times2 =
\qty{700}{m}$. (c) $350/(2\times5) = \qty{35}{m}$ per juta tahun.
\end{solution}

\begin{solution}{exo:b2:oceanic-islands:3}
(a) $5000/17 = \qty{294}{ka}$ di sekitar pulau-pulau utama; $30\,000/70 = \qty{429}{ka}$ di sepanjang
punggungan. (b) Sebagian longsoran terlewat: gunung api yang lebih muda dapat mengubur bekas luka dan puing gunung api yang lebih tua,
dan endapan lama tertutup sedimen. Jumlah sebenarnya lebih besar dan selang sebenarnya
lebih pendek.
\end{solution}

\begin{solution}{exo:b2:oceanic-islands:4}
(a) $4158\times0.3048 = \qty{1267}{m}$; $4610\times0.3048 = \qty{1405}{m}$; $150\times0.3048 =
\qty{46}{m}$. (b) Terumbu pertama tumbuh tidak lebih dalam daripada \qty{46}{m} dan dasarnya kini berada pada
\qty{1267}{m}: gunung api itu telah tenggelam sedikitnya $1267 - 46 = \qty{1221}{m}$ (lebih jika muka laut
saat itu lebih tinggi). (c) Penenggelamannya paling banyak sekitar \qty{1267}{m}, selama sedikitnya \qty{33.9}{Ma}:
paling banyak $1267/33.9 = \qty{37}{m}$ per juta tahun, \qty{0.037}{mm} per tahun.
\end{solution}

\begin{solution}{exo:b2:oceanic-islands:5}
(a) $\dot S = 175/\sqrt t$: \qty{87.5}{m.Ma^{-1}} (\qty{0.0875}{mm.yr^{-1}}) pada \qty{4}{Ma},
\qty{35}{m.Ma^{-1}} (\qty{0.035}{mm.yr^{-1}}) pada \qty{25}{Ma}, \qty{25}{m.Ma^{-1}}
(\qty{0.025}{mm.yr^{-1}}) pada \qty{49}{Ma}. (b) $4.2/0.0875 = 48$ kali. (c) $2.6/0.0875 = 30$
kali: fase pembebanan menenggelamkan pulau jauh lebih cepat daripada pendinginan lempeng muda.
\end{solution}

\begin{solution}{exo:b2:oceanic-islands:6}
(a) $\dot S = (3200/62.8)\,e^{-100/62.8} = 50.96\times0.2035 = \qty{10.37}{m}$ per juta tahun,
\qty{0.0104}{m} per seribu tahun: nilai \qty{0.01}{m} per seribu tahun yang terukur. (b) $e^{-t/62.8}
= 10/50.96$, $t = 62.8\ln5.096 = \qty{102.3}{Ma}$. (c) Hukum akar: $2500 + 350\times5 =
\qty{4250}{m}$; hukum eksponensial: $6400 - 3200\times0.6716 = \qty{4251}{m}$; selisihnya
\qty{1}{m}.
\end{solution}

\begin{solution}{exo:b2:oceanic-islands:7}
(a) $\qty{8.6}{cm.yr^{-1}} = \qty{86}{km.Ma^{-1}}$; $86\times33 = \qty{2838}{km}$. (b) Ya:
\qty{27.7}{Ma} kurang dari \qty{33}{Ma} yang setelahnya terumbu tenggelam, sehingga terumbunya seharusnya
masih hidup. (c) $55\times1.609 = \qty{88.5}{km}$, atau $88.5/86 = \qty{1.0}{Ma}$ lebih tua: sekitar
\qty{28.7}{Ma}. Nilai itu agak kurang dari \qty{33}{Ma}, tetapi \qty{33}{Ma} itu sendiri adalah
angka perkiraan untuk seluruh rantai, dan pertumbuhan juga bergantung pada suhu air di setiap
\omterm{def:b2:oceanic-islands:reefs}{atol}: Kure, \omterm{def:b2:oceanic-islands:reefs}{atol} hidup yang terakhir, memang dekat dengan \omterm{def:b2:oceanic-islands:darwin-point}{titik Darwin}.
\end{solution}

\begin{solution}{exo:b2:oceanic-islands:8}
(a) Kotak: $30\times17\times2 = \qty{1020}{km^3}$; limas: sepertiganya, \qty{340}{km^3}. (b) Satu blok
saja sudah sepersepuluh (kotak) atau sepertiga puluh (limas) seluruh gunung api Kilauea. (c) $23\,000\times
0.1 = \qty{2300}{km^3}$, seperempat Kilauea; taksiran yang sepenuhnya bergantung pada tebal
yang diasumsikan.
\end{solution}

\begin{solution}{exo:b2:oceanic-islands:9}
(a) $1000/2.6 = \qty{385}{ka}$. (b) Ya: \qty{2.6}{m} per seribu tahun sama dengan
\qty{2.6}{mm.yr^{-1}}, kurang dari \qty{4.2}{mm.yr^{-1}}, dengan selisih \qty{1.6}{mm.yr^{-1}}.
(c) Kenaikan relatif $12 + 2.6 = \qty{14.6}{mm.yr^{-1}}$; terumbu tertinggal sebesar $14.6 - 4.2 =
\qty{10.4}{mm.yr^{-1}}$ selama $16.5 - 8.2 = \qty{8.3}{ka}$: puncaknya berakhir pada $10.4\times8.3 =
\qty{86}{m}$. (d) $(12 - 4.2)\times8.3 = \qty{65}{m}$. Dalam kedua kasus terumbu tetap di atas batas pertumbuhan
terdalam tetapi di air yang jauh lebih redup: apakah terumbu itu dapat menyusul bergantung pada bagaimana pertumbuhannya melambat seiring
kedalaman.
\end{solution}

\begin{solution}{exo:b2:oceanic-islands:10}
(a) $S(t_1) = C(\sqrt{t_1 + \Delta} - \sqrt{t_1})$ memiliki turunan $\frac C2(1/\sqrt{t_1 +
\Delta} - 1/\sqrt{t_1}) < 0$: nilainya berkurang seiring $t_1$. (b) $350\sqrt{20} = \qty{1565}{m}$; $350\times
(\sqrt{70} - \sqrt{50}) = 350\times1.2955 = \qty{453}{m}$. (c) Tidak: $\sqrt t$ berhingga di $0$, sehingga
penenggelaman dalam selang mana pun berhingga, \qty{1565}{m} pada (b). Hanya laju pada saat $t = 0$
yang tak berhingga, dan di situ hukum itu tidak lagi dimaksudkan untuk menggambarkan dasar samudra secara rinci.
\end{solution}

\begin{solution}{exo:b2:oceanic-islands:11}
(a) $30\times0.1 = \qty{3}{m}$. (b) Sekitar \qty{120}{m} jika puncaknya berada di muka laut pada
muka laut tinggi: \omterm{def:b2:oceanic-islands:reefs}{atol} itu menjadi pulau batu gamping, hujan melarutkan aragonit karang dan
batuan itu terekristalisasi menjadi kalsit, dengan gua-gua. (c) $120/3 = 40$ siklus, \qty{4}{Ma}. (d) Zona-zona
terlindi adalah masa-masa tersingkap ini, yang tersegel di bawah terumbu yang tumbuh kembali ketika laut kembali.
Selama banyak siklus laut kembali ke muka laut tinggi yang serupa, sehingga $\Delta\eta$ kecil, puncak terumbu
tumbuh kembali sampai muka laut pada setiap muka laut tinggi, dan tebal tudung ditentukan oleh penenggelaman, seperti yang
dinyatakan proposisi itu.
\end{solution}

\begin{solution}{exo:b2:oceanic-islands:12}
(a) $3200\times(e^{-50/62.8} - e^{-120/62.8}) = 3200\times(0.4510 - 0.1480) = \qty{970}{m}$.
(b) $150 + 970 = \qty{1120}{m}$. (c) $700\times1.8288 = \qty{1280}{m}$, \qty{160}{m} lebih dalam daripada
taksiran itu. (d) \omterm{def:b2:oceanic-islands:seamount}{Guyot} itu mungkin tenggelam lebih awal, di atas lempeng yang lebih muda, dan turun lebih jauh; gunung api
tetangga mungkin telah membebaninya; puncaknya mungkin terpangkas lebih rendah oleh gelombang, atau tertutup sedimen;
lempeng di bawahnya mungkin berumur lain; dan muka laut pada waktu tenggelamnya tidak sama dengan
muka laut kini.
\end{solution}

\begin{solution}{pb:b2:oceanic-islands:1}
\textbf{1.} $2500 + 350\sqrt{10} = 2500 + 350\times3.1623 = \qty{3607}{m}$.
\textbf{2.} $3607 + 1000 = \qty{4607}{m}$.
\textbf{3.} $V = \pi r^2h/3 = \pi\times50^2\times4.607/3 = \qty{12061}{km^3}$.
\textbf{4.} $0.85\times12\,061 = \qty{10252}{km^3}$ dalam \qty{1}{Ma}: \qty{0.0103}{km^3} per tahun,
sepuluh kali lebih kecil daripada Kilauea kini.
\textbf{5.} $95\,600/12\,061 = 7.9$: sekitar delapan.
\textbf{6.} Lava toleiitik, dari \omterm{def:b2:oceanic-islands:stages}{tahap perisai}; lava alkali tahap pascaperisai menudunginya.
\textbf{7.} $\qty{1}{km}/\qty{1}{Ma} = \qty{1}{m}$ per seribu tahun, 2.6 kali lebih lambat daripada di lepas
Pulau Hawaii.
\textbf{8.} Ya: $4.2 - 1.0 = \qty{3.2}{mm.yr^{-1}}$ tersisa.
\textbf{9.} $12 + 1 - 4.2 = \qty{8.8}{mm}$ per tahun.
\textbf{10.} $150/8.8 = \qty{17}{ka}$. Sebuah teras terumbu tenggelam di lereng, anak tangga dari jenis yang
ditemukan di lepas Pulau Hawaii.
\textbf{11.} $350\times(\sqrt{22} - \sqrt{10}) = 350\times(4.6904 - 3.1623) = \qty{535}{m}$.
\textbf{12.} $1000 - 535 = \qty{465}{m}$, dengan laju $465/12 = \qty{38.8}{m}$ per juta tahun.
\textbf{13.} $350/(2\sqrt{10}) = \qty{55.3}{m.Ma^{-1}}$ pada awalnya; $350/(2\sqrt{22}) =
\qty{37.3}{m.Ma^{-1}}$ pada akhirnya.
\textbf{14.} $2500 + 350\times4.6904 = \qty{4142}{m}$.
\textbf{15.} Pulau-pulau Hawaii yang besar menjulang sampai sekitar empat kali lebih tinggi, sehingga ada jauh lebih banyak
batuan di atas muka laut yang harus disingkirkan; \omterm{def:b2:oceanic-islands:stages}{tahap perisainya} juga berlangsung lebih lama.
\textbf{16.} Gelombang, yang memundurkan pantai dan memangkas dataran abrasi; longsoran dan \omterm{def:b2:oceanic-islands:collapse}{keruntuhan
lereng}.
\textbf{17.} \omterm{def:b2:oceanic-islands:reefs}{Terumbu penghalang}, dengan laguna di antara terumbu dan sisa terakhir pantai vulkanik.
\textbf{18.} $350\times(\sqrt{70} - \sqrt{22}) = 350\times(8.3666 - 4.6904) = \qty{1287}{m}$.
\textbf{19.} $1287/48 = \qty{26.8}{m}$ per juta tahun, \qty{0.027}{mm} per tahun: $4.2/0.0268 =
157$ kali lebih kecil daripada laju yang dapat dicapai pertumbuhan terumbu.
\textbf{20.} \qty{20}{m} lebih tebal daripada batu gamping di atas basal di bawah Pulau Parry
(\qty{1267}{m}), dan \qty{118}{m} lebih tipis daripada di bawah Elugelab (\qty{1405}{m}): ukurannya sama.
\textbf{21.} $2500 + 350\times8.3666 = \qty{5428}{m}$.
\textbf{22.} $3200\times(e^{-70/62.8} - e^{-120/62.8}) = 3200\times(0.3280 - 0.1480) =
\qty{576}{m}$.
\textbf{23.} $150 + 576 = \qty{726}{m}$.
\textbf{24.} Tudung terumbu \omterm{def:b2:oceanic-islands:reefs}{atol} Petrel setebal \textbf{\qty{1287}{m}}, sekitar \qty{1.3}{km}
batu gamping terumbu, seperti di bawah Enewetak.
\end{solution}
